%! Exponential Growth and Decay — Unit 6, Lesson 6.1 — Homework (Answer Key)
\documentclass[10pt]{article}
\usepackage{algebra2-article}
\usepackage{algebra2-key}   % pulls in -boxes; do NOT also load -boxes
\newcommand{\TallMath}[1]{$\displaystyle #1\rule[-1.4em]{0pt}{3.2em}$}
\newcommand{\lgrid}[4]{%
  \draw[step=1, linegray!40, very thin] (#1,#3) grid (#2,#4);
  \draw[->, semithick, charcoal] (#1,0)--(#2,0) node[below right, font=\scriptsize]{$x$};
  \draw[->, semithick, charcoal] (0,#3)--(0,#4) node[left, font=\scriptsize]{$y$};}
% #2 is ln(b):  ln2 = 0.693147   ln3 = 1.098612
\newcommand{\expcurve}[4]{\draw[very thick, forest, domain=#3:#4, samples=120, smooth]
  plot (\x,{#1*exp((#2)*(\x))});}

\begin{document}

\pageheader{Unit 6, Lesson 6.1}{Homework --- Answer Key}
\namedateperiod

\vspace{0.10in}

\begin{notesbox}{Practice}
Show the two numbers ($a$ and $b$) for every model. Round money to the nearest cent and counts of
living things down to a whole number.

\begin{enumerate}[leftmargin=1.9em, itemsep=12pt]
  \item \textbf{Percent and factor, both directions.}

    \textbf{(a)} Write the factor $b$ for each description.
    \begin{center}
    \renewcommand{\arraystretch}{1.5}
    \begin{tabularx}{\linewidth}{@{}X l l X@{}}
    \hline
    \textbf{The sentence says\dots} & \textbf{Rate $r$} & \textbf{Grows or shrinks?} & \textbf{Factor $b$} \\
    \hline
    grows by $12\%$ each year & \ans{0.12} & \ans{grows} & \ans{$1.12$} \\
    decreases by $5\%$ each month & \ans{0.05} & \ans{shrinks} & \ans{$0.95$} \\
    loses $8\%$ of it each hour & \ans{0.08} & \ans{shrinks} & \ans{$0.92$} \\
    doubles each day & \ans{$1$ ($100\%$)} & \ans{grows} & \ans{$2$} \\
    increases by $0.75\%$ each year & \ans{0.0075} & \ans{grows} & \ans{$1.0075$} \\
    \hline
    \end{tabularx}
    \end{center}

    \textbf{(b)} Now go backwards. Each model is built --- describe it in words.
    \begin{center}
    \renewcommand{\arraystretch}{1.5}
    \begin{tabularx}{\linewidth}{@{}l X X X@{}}
    \hline
    \textbf{Model} & \textbf{Initial value} & \textbf{Percent rate} & \textbf{Increase or decrease?} \\
    \hline
    $y=300(1.09)^{t}$ & \ans{300} & \ans{$9\%$} & \ans{increase} \\
    $y=64(0.85)^{t}$ & \ans{64} & \ans{$15\%$} & \ans{decrease} \\
    $y=1500(1.5)^{t}$ & \ans{1500} & \ans{$50\%$} & \ans{increase} \\
    \hline
    \end{tabularx}
    \end{center}

  \item \textbf{From tables.} For each table, decide whether it can be modeled by $y=ab^{x}$. If it
    can, give $a$, $b$, and the equation. If it cannot, say what family it belongs to and why.
    \begin{center}
    \renewcommand{\arraystretch}{1.25}
    \begin{tabular}{|c|c|c|c|c|c|}
    \hline
    $x$ & $0$ & $1$ & $2$ & $3$ & $4$ \\ \hline
    A: $y$ & $6$ & $18$ & $54$ & $162$ & $486$ \\ \hline
    B: $y$ & $160$ & $80$ & $40$ & $20$ & $10$ \\ \hline
    C: $y$ & $2$ & $5$ & $8$ & $11$ & $14$ \\ \hline
    \end{tabular}
    \end{center}
    A: \; $a=$ \ans{6} \quad $b=$ \ans{3} \quad $y=$ \ans{$6\cdot 3^{x}$} \\[3pt]
    B: \; $a=$ \ans{160} \quad $b=$ \ans{$\frac12$} \quad $y=$ \ans{$160\left(\frac12\right)^{x}$} \\[3pt]
    C: \; \ans{Not exponential --- \emph{linear}. The ratios ($\frac52$, $\frac85$, $\frac{11}{8}$) are not constant, but the differences are a constant $+3$, so $y=3x+2$.}

  \item \textbf{From a graph.} The curve below is $y=ab^{x}$, with two points labeled.
    \begin{center}
    \begin{minipage}[c]{0.42\linewidth}
    \centering
    \begin{tikzpicture}[scale=0.42]
      \lgrid{-1}{6}{-1}{10}
      \begin{scope}
        \clip (-1,-1) rectangle (6,10);
        \expcurve{8}{-0.693147}{-0.32}{6}
      \end{scope}
      \draw[navy, dashed, semithick] (-1,0)--(6,0);
      \fill[charcoal] (0,8) circle (3pt) node[right, font=\scriptsize]{$(0,8)$};
      \fill[charcoal] (3,1) circle (3pt) node[above right, font=\scriptsize]{$(3,1)$};
    \end{tikzpicture}
    \end{minipage}
    \begin{minipage}[c]{0.54\linewidth}
    $a=$ \ans{8} \; (read it off the $y$-axis) \\[6pt]
    Substituting $(3,1)$: \; $1=a\,b^{3}$, so $b^{3}=$ \ans{$\frac18$} \\[6pt]
    and $b=$ \ans{$\frac12$}. \quad Growth or decay? \ans{decay} \\[6pt]
    $y=$ \ans{$8\left(\frac12\right)^{x}$} \\[6pt]
    Range: \ans{$(0,\infty)$} \quad Asymptote: \ans{$y=0$}
    \end{minipage}
    \end{center}

  \item \textbf{From stories.} Build each model, evaluate it, and interpret the result.

    \textbf{(a)} A culture starts with $250$ bacteria cells and the count increases by $40\%$ every
    hour. \\[3pt]
    $C(t)=$ \ans{$250(1.4)^{t}$} \quad $C(6)=$ \ans{$250(7.5295)$} $\approx$ \ans{$1882$} cells \\[3pt]
    What does $C(6)$ mean? \ansline{Six hours after the culture was started there are about $1882$
    cells --- the colony has grown to roughly seven and a half times its original size.}

    \vspace{4pt}
    \textbf{(b)} A phone bought for \$$900$ loses $22\%$ of its value each year. \\[3pt]
    $V(t)=$ \ans{$900(0.78)^{t}$} \quad $V(4)=$ \ans{$900(0.37015)$} $\approx$ \ans{\$$333.14$} \\[3pt]
    What does $V(4)$ mean? \ansline{Four years after it was bought the phone is worth about
    \$$333.14$, a little over a third of what it cost.}

  \item \textbf{Error analysis.} For each student, name the error, say what their model
    \emph{actually} describes, and write the correct model.

    \textbf{(a)} ``A \$$18{,}000$ truck loses $15\%$ of its value each year.'' \\
    \hspace*{1.2em} Jonah wrote $V(t) = 18{,}000(1.15)^{t}$.
    \ansline{Jonah \emph{added} the rate instead of subtracting it. His base is greater than $1$, so
    his model has the truck \emph{gaining} $15\%$ a year --- worth more than \$$36{,}000$ after five
    years. Losing $15\%$ leaves $85\%$, so $b=1-0.15=0.85$ and
    $V(t)=18{,}000(0.85)^{t}$.}

    \textbf{(b)} ``A colony of $250$ bees grows $7\%$ each week.'' \\
    \hspace*{1.2em} Priya wrote $P(w) = 250(0.07)^{w}$.
    \ansline{Priya used the \emph{percent} as the base. Her model keeps only $7\%$ of the colony each
    week --- a $93\%$ collapse, down to about $17$ bees after one week. Growing $7\%$ means having
    $107\%$, so $b=1+0.07=1.07$ and $P(w)=250(1.07)^{w}$.}
\end{enumerate}
\end{notesbox}

\begin{extensionbox}
\textbf{Two towns, ten years.} Ashford has $20{,}000$ people and is \textbf{losing $2\%$} of its
population each year. Brookvale has $12{,}000$ people and is \textbf{gaining $5\%$} each year.
\begin{center}
\renewcommand{\arraystretch}{1.25}
\begin{tabular}{|c|c|c|c|c|c|c|}
\hline
Year $t$ & $0$ & $2$ & $4$ & $6$ & $8$ & $10$ \\ \hline
Ashford & $20{,}000$ & $19{,}208$ & $18{,}447$ & $17{,}717$ & $17{,}015$ & $16{,}341$ \\ \hline
Brookvale & $12{,}000$ & $13{,}230$ & $14{,}586$ & $16{,}081$ & $17{,}730$ & $19{,}547$ \\ \hline
\end{tabular}
\end{center}
\begin{enumerate}[label=(\alph*), leftmargin=1.8em, itemsep=5pt]
  \item Write both models, $A(t)$ and $B(t)$.
  \item Verify one entry in each row of the table by evaluating your model.
  \item Brookvale starts $8{,}000$ people behind. Between which two years in the table does it pass
        Ashford? What tells you that from the table alone?
  \item To find the crossing \emph{exactly} you would have to solve
        $20{,}000(0.98)^{t} = 12{,}000(1.05)^{t}$. Try rewriting both sides with a common base ---
        what goes wrong?
  \item Ashford's model has horizontal asymptote $A=0$. What does that claim about the town, and how
        far into the future would you trust this model? Explain.
\end{enumerate}
\ansline{(a) $A(t)=20{,}000(0.98)^{t}$ \; and \; $B(t)=12{,}000(1.05)^{t}$. Losing $2\%$ leaves
$98\%$; gaining $5\%$ gives $105\%$.\\[3pt]
(b) For example $A(4)=20{,}000(0.98)^{4}=20{,}000(0.92237)\approx 18{,}447$ and
$B(4)=12{,}000(1.05)^{4}=12{,}000(1.21551)\approx 14{,}586$ --- both match.\\[3pt]
(c) Between year $6$ and year $8$. At $t=6$ Ashford is still ahead ($17{,}717$ vs.\ $16{,}081$); at
$t=8$ Brookvale is ahead ($17{,}730$ vs.\ $17{,}015$), so the two populations are equal somewhere in
between. (Checking year by year puts it a little past $t=7$.)\\[3pt]
(d) There is no common base: $0.98$ and $1.05$ are not powers of any one number you can write down,
so ``rewrite both sides over the same base and equate the exponents'' has nothing to work with. You
can \emph{read} the crossing off a graph or narrow it down in a table, but solving it exactly needs a
tool we do not have yet --- that is what Unit 7 is for.\\[3pt]
(e) It claims Ashford never fully empties: each year keeps $98\%$ of a positive number, which stays
positive forever, so the model has no $t$-intercept. That is fine for the next decade or two, when
$2\%$ is a plausible steady rate. It is not believable across a century --- populations do not decay
smoothly forever, and a town of $3$ people is not a town. Trust the model in the range the data came
from, not in the tail.}
\end{extensionbox}

\begin{spiralbox}
\textbf{Coming next --- Lesson 6.2: Graphing Exponential Functions \& Transformations.} You can now
build $y=ab^{x}$. Next you will move it: shift it, stretch it, and flip it, using the same
transformation rules you have applied since Unit 1. One of them behaves differently here than
anywhere else in the course --- adding $k$ to the outside, $f(x)+k$, is the only transformation that
\emph{moves the horizontal asymptote}, and therefore the only one that changes the range. Also worth
noticing today: $y=8\left(\frac12\right)^{x}$ in problem 3 could just as well have been written
$y=8\cdot 2^{-x}$. In 6.2 that stops being a curiosity and becomes a rule --- reflecting an
exponential across the $y$-axis turns growth into decay.
\end{spiralbox}

\begin{teachernote}
\textbf{Problem 1(b) is the item to grade first.} Running the translation backwards is what proves a
student understands the rule rather than pattern-matching it, and the $0.85 \Rightarrow 15\%$
\emph{decrease} row is where the sign error shows up. The $1.5 \Rightarrow 50\%$ row catches students
who read the base as the percent.

\textbf{Problem 2, Table C} is the deliberate trap: it is the only table that is not exponential, and
it must be diagnosed with the Lesson 6.0 fingerprint test rather than by eye. Accept ``linear'' with
the constant difference cited; the equation $y=3x+2$ is a bonus.

\textbf{Problem 3} is Section 4's Case 1 with a fractional answer --- $b^{3}=\frac18$ gives
$b=\frac12$, not $b=\frac13$. Watch for students dividing the exponent instead of taking a root.

\textbf{Problem 5 carries the lesson.} Jonah and Priya make the \emph{two different} errors, and the
assignment asks students to name each one, not just fix it. A complete answer says which side of $1$
the wrong base fell on and what percent change it really encodes. If a student can write both
paragraphs, they own the translation.

\textbf{The Extension sets up two future lessons.} Part (d) is the wall Lesson 6.4 ends on, arriving
a few days early on purpose --- do not let students grind at it, and do not hint that logarithms
exist unless someone asks. Part (e) reruns Lesson 6.0's asymptote-believability question in a new
context, and part (c) is the reading-a-crossing-from-a-table skill that 6.3 will make algebraic in
the special case where a common base \emph{does} exist.
\end{teachernote}

\end{document}
